\documentclass[10pt,a4paper]{amsart}
\usepackage[margin=19mm]{geometry}
\usepackage{amsmath,amssymb,mathtools}
\usepackage[scr=rsfs,cal=euler]{mathalfa}
\usepackage{xfrac}
\usepackage[unicode,hidelinks,bookmarks=false]{hyperref}
\newcommand{\ie}{i.e.\ }
\newcommand{\cf}{cf.\ }
\newcommand{\resp}{respectively\ }
\newcommand{\Subsection}{\subsection}
\providecommand{\nobreakdash}{\nobreak\hbox{-}}
\setlength{\emergencystretch}{2em}
\multlinegap0pt
\usepackage{amssymb}
\usepackage{mathtools}
\usepackage{bm}
\usepackage[normalem]{ulem}
\usepackage{csquotes}
\newtheorem{lemma}{Lemma}
\newcommand{\C}{\mathbb{C}}
\DeclareMathOperator{\spec}{spec}
\title{Random Schr\"odinger operators on manifolds and abstract bounds for multiplier-type operators}
\author{Jean-Claude Cuenin, Konstantin Merz, Eduard Stefanescu}
\date{Source: arXiv:2606.19075v2 (CC BY 4.0)}
\begin{document}
\maketitle

\noindent\emph{Source and licence.} Random Schr\"odinger operators on manifolds and abstract bounds for multiplier-type operators. Version arXiv:2606.19075v2 (CC BY 4.0), distributed by arXiv under the Creative Commons Attribution 4.0 International licence, http://creativecommons.org/licenses/by/4.0/, as recorded in the arXiv licence metadata for this exact version and shown on the abstract page https://arxiv.org/abs/2606.19075v2. Full licence notice: \texttt{licenses/arxiv-2606.19075-CC-BY-4.0.txt}.\par\smallskip

\noindent\textit{Editorial scope.} Counted unit: the article's lemma lem:BS-perturb-z (source0.tex lines 1591--1597). The article's complete original proof of it is delivered with it (source0.tex lines 1599--1701, one \texttt{proof} environment of 103 lines). No mathematics is written, repaired or re-derived: the statement and the proof are the article's own lines, in the article's own notation, and no retained line is substituted or re-flowed.  No further source-local context is needed: every cross-reference of every retained line resolves inside the two delivered spans.  Nothing was omitted from any retained span, so the package carries no omission record.  No retained line contains a citation, so the wrapper carries no bibliography and no bibliography entry.  No retained line contains a figure environment, an image-inclusion command or drawing code, so no image is delivered and the package declares no asset dependency.  Editorial formatting only, in addition: an AMS \texttt{amsart} preamble that restates the article's own declarations and macros used by the retained lines, namely the environments \texttt{lemma} on separate counters, together with the standard packages those lines need.  The retained environments print in this document as Lemma 1 in order of appearance, whereas the article prints the same items with the numbering of its own full document.  The complete native source of the article as distributed in the arXiv e-print for this exact version is bundled unchanged as \texttt{sources/203114/source0.tex}.
\begin{lemma}\label{lem:BS-perturb-z}
  For any \(z,z'\in \C\setminus \spec(-\Delta_g)\) with $d(z),d(z')\geq  \rho>0$, we have
  \[
    \|\mathcal{K}_{\omega,\lambda}(z)-\mathcal{K}_{\omega,\lambda}(z')\|
    \lesssim |z-z'| \, \rho^{-2}\, \|V_{\omega,\lambda}\|_{L^{\infty}(M)}.
  \]
\end{lemma}

\begin{proof}
  Set
\[
R(w):=(-\Delta_g-w)^{-1/2},
\qquad
w\in \C\setminus \spec(-\Delta_g).
\]
Then
\[
\mathcal{K}_{\omega,\lambda}(w)=|R(w)|V_{\omega,\lambda}R(w),
\]
and hence
\[
\mathcal{K}_{\omega,\lambda}(z)-\mathcal{K}_{\omega,\lambda}(z')
=
\bigl(|R(z)|-|R(z')|\bigr)V_{\omega,\lambda}R(z)
+
|R(z')|V_{\omega,\lambda}\bigl(R(z)-R(z')\bigr).
\]
Therefore, taking the operator norm,
\begin{align*}
\|\mathcal{K}_{\omega,\lambda}(z)-\mathcal{K}_{\omega,\lambda}(z')\|
&\leq
\|R(z)-R(z')\|\,\|V_{\omega,\lambda}\|\,\|R(z)\| \\
&\qquad
+\|R(z')\|\,\|V_{\omega,\lambda}\|\,\|R(z)-R(z')\|,
\end{align*}
where $\|V_{\omega,\lambda}\|=\|V_{\omega,\lambda}\|_{L^{\infty}(M)}$.
It remains to estimate \(\|R(w)\|\) and \(\|R(z)-R(z')\|\).

By the spectral theorem,
\begin{align}\label{eq:trivial bound R(w)}
\|R(w)\|
=
\sup_{\mu\in \spec(-\Delta_g)} |\mu-w|^{-1/2}
=
d(w)^{-1/2}.
\end{align}

For the difference, define for \(\mu\in \spec(-\Delta_g)\)
\[
f_w(\mu):=|\mu-w|^{-1/2}.
\]
Then
\[
\|R(z)-R(z')\|
=
\sup_{\mu\in \spec(-\Delta_g)} |f_z(\mu)-f_{z'}(\mu)|.
\]
Using 
\begin{align*}
\bigl||a|^{-1/2}-|b|^{-1/2}\bigr|
&=
\frac{\bigl||b|^{1/2}-|a|^{1/2}\bigr|}{|a|^{1/2}|b|^{1/2}}
=
\frac{|\,|b|-|a|\,|}{|a|^{1/2}|b|^{1/2}(|a|^{1/2}+|b|^{1/2})}\\
&\leq
\frac{|a-b|}{|a|^{1/2}|b|^{1/2}(|a|^{1/2}+|b|^{1/2})},
\end{align*}
with \(a=\mu-z\) and \(b=\mu-z'\), we obtain
\[
|f_z(\mu)-f_{z'}(\mu)|
\leq
\frac{|z-z'|}{|\mu-z|^{1/2}|\mu-z'|^{1/2}
\bigl(|\mu-z|^{1/2}+|\mu-z'|^{1/2}\bigr)}.
\]
Since
\[
\frac{1}{x^{1/2}y^{1/2}(x^{1/2}+y^{1/2})}
\leq
\min(x^{-1/2}y^{-1},x^{-1}y^{-1/2})
\leq
x^{-3/2}+y^{-3/2},
\]
for all \(x,y>0\), it follows that
\[
|f_z(\mu)-f_{z'}(\mu)|
\lesssim
|z-z'|\bigl(|\mu-z|^{-3/2}+|\mu-z'|^{-3/2}\bigr).
\]
Taking the supremum over \(\mu\in \spec(-\Delta_g)\) yields
\[
\|R(z)-R(z')\|
\lesssim
|z-z'|\bigl(d(z)^{-3/2}+d(z')^{-3/2}\bigr).
\]

Substituting these bounds into the previous estimate gives
\begin{align*}
\|\mathcal{K}_{\omega,\lambda}(z)-\mathcal{K}_{\omega,\lambda}(z')\|
&\lesssim
|z-z'|\,\|V_{\omega,\lambda}\|
\bigl(d(z)^{-3/2}+d(z')^{-3/2}\bigr)
\bigl(d(z)^{-1/2}+d(z')^{-1/2}\bigr).
\end{align*}
If \(d(z),d(z')\geq \rho\), then
\[
\|\mathcal{K}_{\omega,\lambda}(z)-\mathcal{K}_{\omega,\lambda}(z')\|
\lesssim
|z-z'|\,\rho^{-2}\,\|V_{\omega,\lambda}\|_{L^{\infty}(M)}.
\]
This completes the proof.
\end{proof}

\end{document}
